\documentclass[11pt,a4paper]{article}

\usepackage[margin=2.0cm]{geometry}
\usepackage[T1]{fontenc}
\usepackage[utf8]{inputenc}
\usepackage{titlesec}
\usepackage{xcolor}
\usepackage{fancyhdr}
\usepackage{hyperref}
\usepackage{parskip}
\usepackage{fancyvrb}
\usepackage{tcolorbox}

\definecolor{accentblue}{rgb}{0.12,0.31,0.57}

\pagestyle{fancy}
\fancyhf{}
\lhead{\textbf{Function Syntax Reference}}
\rhead{Environmental Monitor Project}
\cfoot{Page \thepage}
\renewcommand{\headrulewidth}{0.4pt}

\titleformat{\section}{\Large\bfseries\color{accentblue}}{}{0em}{}[\vspace{-0.5em}\rule{\textwidth}{0.4pt}]
\titleformat{\subsection}{\large\bfseries\color{accentblue}}{}{0em}{}

\newtcolorbox{infobox}[1][]{
  colback=blue!5,
  colframe=accentblue,
  fonttitle=\bfseries,
  title={#1},
  arc=2pt,
  boxrule=0.5pt
}

\begin{document}

\begin{center}
  \vspace*{1.2cm}
  {\Huge\bfseries\color{accentblue} Function Syntax Reference}\\[1em]
  {\Large Environmental Monitor Actuator and Sender Nodes}\\[1em]
  {\large LaTeX wrapper for the full syntax guide}\\[1.5em]
  \rule{0.8\textwidth}{0.4pt}
\end{center}

\section{Overview}

This document is the LaTeX version of your function syntax guide. It includes the full text of the reference so you can compile it as a standalone PDF for study and interview preparation.

\begin{infobox}[What This Document Covers]
This reference explains:
\begin{itemize}
  \item GPIO, NVIC, CAN, FreeRTOS, HAL, Timer, and RCC function syntax
  \item Parameter meanings and return values
  \item Real examples taken from your codebase
  \item Plain-English explanations of what each call does in your project
\end{itemize}
\end{infobox}

\section{Full Reference}

{\footnotesize
\VerbatimInput[
  frame=single,
  framesep=6pt,
  breaklines=true,
  breakanywhere=true,
  fontsize=\footnotesize,
  commandchars=\\\{\}
]{Function_Syntax_Reference.txt}
}

\end{document}
